\documentclass[11pt]{article}
\usepackage[T1]{fontenc}
\usepackage[english]{babel}
\usepackage{amsmath,amssymb,amsthm}
\usepackage{hyperref}

\title{On Large Deviations in the ABC Conjecture and Prime Factor Constraints}
\author{Thomas Hoffbauer}
\date{}

\newtheorem{theorem}{Theorem}
\newtheorem{lemma}{Lemma}
\newtheorem{proposition}{Proposition}
\newtheorem{remark}{Remark}

\begin{document}
	\maketitle
	
	\begin{abstract}
		We study large deviations in the ABC conjecture via the quantity
		[
		\Delta(a,b,c) := \log c - \log \mathrm{rad}(abc).
		]
		Numerical experiments indicate a strong negative correlation between $\Delta$
		and the number $\omega(abc)$ of distinct prime factors.
		
		We formulate a structural hypothesis stating that large values of $\Delta$
		can only occur when the prime support is small, and provide a partial
		theoretical explanation based on S-unit equations and bounds from
		Baker theory.
	\end{abstract}
	
	\section{Introduction}
	
	The ABC conjecture states that for coprime integers $a,b$ with $a+b=c$,
	and for every $\varepsilon>0$, there exists $K_\varepsilon$ such that
	[
	c \le K_\varepsilon \cdot \mathrm{rad}(abc)^{1+\varepsilon}.
	]
	
	We define
	[
	\Delta(a,b,c) := \log c - \log \mathrm{rad}(abc).
	]
	
	The conjecture implies that $\Delta$ cannot grow significantly.
	
	\section{Numerical Evidence}
	
	Computations for coprime pairs $(a,b)$ with $a,b \le N$ (up to several million)
	show:
	
	\begin{itemize}
		\item $\Delta$ is typically negative,
		\item large values of $\Delta$ are rare,
		\item large $\Delta$ occurs only when $\omega(abc)$ is small.
	\end{itemize}
	
	\subsection{Structural observation}
	
	In all observed cases with large $\Delta$:
	
	\begin{itemize}
		\item the number of distinct prime factors is small,
		\item high prime powers occur,
		\item $c$ is often close to a prime power.
	\end{itemize}
	
	These observations motivate the following hypothesis.
	
	\section{Structural Constraint}
	
	\begin{lemma}[Structural Constraint]
		For every $\delta>0$, there exists a constant $K(\delta)$ such that
		[
		\Delta(a,b,c) \ge \delta ;\Rightarrow; \omega(abc) \le K(\delta).
		]
	\end{lemma}
	
	\begin{remark}
		This lemma is not proven here. It is supported by numerical evidence
		and consistency with known diophantine constraints.
	\end{remark}
	
	\section{Connection to S-Unit Equations}
	
	Let
	[
	S := {p : p \mid abc}, \quad s = |S|.
	]
	
	Define
	[
	x = \frac{a}{c}, \quad y = \frac{b}{c}.
	]
	Then
	[
	x + y = 1,
	]
	where $x,y$ are S-units.
	
	By classical results (see \cite{EvertseGyory}),
	the number of such solutions is finite for fixed $S$.
	
	\section{Exponent Constraints from Baker Theory}
	
	We write
	[
	c = \prod_{p \in S} p^{\gamma_p}.
	]
	
	Bounds from Baker theory imply that the exponents $\gamma_p$ cannot grow
	arbitrarily when the size $s = |S|$ increases.
	
	More precisely, there exists a function $E(s)$ with at most exponential growth
	such that
	[
	\max_{p \in S} \gamma_p \le E(s).
	]
	
	\begin{proposition}[Heuristic exponent bound]
		For solutions of $a+b=c$ with prime support $S$, one has
		[
		\Delta(a,b,c) \le (E(s)-1)\sum_{p \in S} \log p.
		]
	\end{proposition}
	
	\begin{proof}[Heuristic argument]
		We write
		[
		\Delta = \sum_{p \in S} (\gamma_p - 1)\log p.
		]
		If $s$ is large, Baker-type bounds imply that most exponents $\gamma_p$
		must remain small, preventing large contributions to $\Delta$.
	\end{proof}
	
	\section{Heuristic Consequences}
	
	Using the prime number theorem,
	[
	\sum_{p \in S} \log p \gtrsim s \log s.
	]
	
	Thus,
	[
	\Delta \lesssim E(s), s \log s.
	]
	
	Since $E(s)$ grows at most exponentially, this suggests that large $\Delta$
	can only occur when $s$ is small.
	
	\section{Quantitative Behaviour}
	
	The above argument suggests the existence of a monotone function $K(\delta)$ such that
	[
	\Delta \ge \delta ;\Rightarrow; \omega(abc) \le K(\delta),
	]
	with $K(\delta)$ growing slowly as $\delta \to 0$.
	
	\section{Discussion}
	
	This perspective suggests that large deviations in the ABC conjecture
	correspond to highly constrained configurations of prime factors.
	
	In particular, large $\Delta$ appears to require strong concentration
	of multiplicative structure in a small number of primes.
	
	\section{Outlook}
	
	A possible route toward the ABC conjecture is to rigorously establish
	the structural constraint or to derive it from refined diophantine bounds.
	
	\begin{thebibliography}{9}
		
		\bibitem{BakerWustholz}
		A.~Baker and G.~Wüstholz,
		\textit{Logarithmic Forms and Diophantine Geometry},
		Cambridge University Press, 2007.
		
		\bibitem{EvertseGyory}
		J.-H.~Evertse and K.~Győry,
		\textit{Unit Equations in Diophantine Number Theory},
		Cambridge University Press, 2015.
		
	\end{thebibliography}
	
\end{document}
